\chapter{Por que vigiar um poço de petróleo}
\label{cap:introducao}

\begin{objetivos}
Ao final deste capítulo, você será capaz de:
\begin{itemize}[leftmargin=*,itemsep=0.2em]
  \item explicar por que a operação de poços submarinos exige monitoramento contínuo e
        automatizado, em termos de segurança, integridade e economia;
  \item quantificar o impacto do tempo não produtivo (NPT) e relacioná-lo à detecção
        precoce de anomalias;
  \item distinguir os três problemas computacionais envolvidos --- detecção,
        classificação e explicação --- e enunciar por que eles não são o mesmo problema;
  \item enunciar a hipótese de mundo fechado e reconhecer suas manifestações práticas;
  \item descrever o percurso argumentativo do livro e escolher a trilha de leitura
        adequada ao seu objetivo.
\end{itemize}
\end{objetivos}

\section{Um poço é um sistema que você não pode ver}

A engenharia costuma lidar com sistemas que podem ser inspecionados. Uma ponte pode
ser percorrida a pé; um vaso de pressão pode ser ensaiado por ultrassom; uma turbina
pode ser aberta. Um poço submarino de petróleo em lâmina d'água de dois mil metros não
pode ser nada disso. Entre o operador na sala de controle e o reservatório a seis mil
metros abaixo da linha d'água existem, no melhor dos casos, algumas dezenas de canais
de telemetria --- pressões, temperaturas, vazões --- amostrados a cada poucos segundos.

Tudo o que se sabe sobre o estado do poço vem desses canais. Se um deles congela, se
um transdutor deriva, se um evento se manifesta em uma escala temporal mais lenta do
que a janela de observação do operador, a observação pode ficar ausente, enviesada ou
insuficiente para identificar o estado. A
consequência é direta: \emph{o monitoramento de poços é, antes de tudo, um problema de
inferência a partir de observação parcial}, uma das classes de problema em que o
aprendizado de máquina pode ser útil.

\subsection{A escala do problema}

Uma única unidade flutuante de produção, armazenamento e transferência (FPSO) pode
estar conectada a algumas dezenas de poços produtores e injetores. Uma operadora de
porte médio opera dezenas de FPSOs. O sistema descrito no \cref{cap:dual} deste livro
foi implantado, à época de sua avaliação, em mais de vinte unidades e mais de duzentos
e cinquenta poços submarinos nas bacias de Santos e de Campos.

Multiplique: se cada poço fornece cinco sinais amostrados a cada segundo, duzentos e
cinquenta poços produzem mais de cem milhões de leituras por dia. Nenhuma equipe de
operação lê cem milhões de números por dia. O que se lê, na prática, são alarmes --- e
alarmes são decisões automatizadas. A qualidade do monitoramento é, portanto, a
qualidade dos algoritmos que decidem quando tocar um alarme.

\subsection{O que dá errado}
\label{sec:introducao-oquedaerrado}

Os eventos indesejáveis que acometem poços de petróleo formam uma taxonomia razoavelmente
estável, detalhada no \cref{cap:dominio}. Por ora, quatro exemplos bastam para dar
concretude ao problema:

\begin{itemize}[leftmargin=*,itemsep=0.4em]
  \item \textbf{Formação de hidratos.} Em águas profundas, a combinação de alta pressão
        e baixa temperatura leva moléculas de água e gás leve a formar um sólido
        cristalino semelhante a gelo. Um tampão de hidrato pode obstruir completamente
        uma linha de produção. A remoção exige despressurização controlada, injeção de
        inibidor ou intervenção submarina, com custo que se mede em milhões de dólares e
        prazo que se mede em semanas.
  \item \textbf{Fechamento espúrio da válvula de segurança de subsuperfície (DHSV).}
        A DHSV é o dispositivo que fecha o poço em situação de emergência. Quando ela
        fecha sem que haja emergência --- por perda de pressão hidráulica na linha de
        controle, por exemplo --- a produção é interrompida. Não é um evento perigoso;
        é um evento caro, e sua reabertura exige diagnóstico correto.
  \item \textbf{Golfadas severas.} O escoamento multifásico em um \ing{riser} vertical
        pode entrar em regime oscilatório de grande amplitude, em que golfadas de
        líquido e bolhas de gás se alternam. O resultado são oscilações de pressão que
        castigam equipamentos de superfície e degradam a eficiência de separação.
  \item \textbf{Incrustação na válvula reguladora de produção (PCK).} Depósitos
        minerais reduzem gradualmente a área de passagem. O sintoma é lento, cumulativo
        e facilmente confundido com declínio natural do reservatório.
\end{itemize}

Observe a heterogeneidade. O hidrato é um evento de horas; a incrustação, de meses. O
fechamento de DHSV é uma transição quase instantânea; a golfada é um regime
estacionário oscilatório. Um método que detecte bem um deles não detectará
necessariamente bem os outros --- e essa observação, que parece trivial aqui, será a
justificativa técnica para praticamente todas as escolhas de arquitetura discutidas
neste livro.

\section{O custo de não saber}

\begin{definicao}{Tempo não produtivo}
Chama-se \termo{tempo não produtivo}{non-productive time} (NPT) o intervalo em que um
poço, sonda ou instalação está indisponível para sua função por causa não planejada.
O NPT é contabilizado em horas ou dias e convertido em custo pela soma do custo de
oportunidade da produção perdida com o custo direto da intervenção corretiva.
\end{definicao}

O NPT é a métrica econômica que conecta a estatística de detecção à decisão de
investimento. Um sistema de monitoramento que antecipa a identificação de um evento em
algumas horas pode transformar uma intervenção submarina em um ajuste de \ing{choke}
feito da sala de controle.

\begin{exemplo}{Estimativa de NPT em um cenário hipotético}
Considere um poço que produziria $q=\num{10000}$ barris por dia, preço de referência
$p=\num{70}$ dólares por barril, doze dias de NPT e quatro dias de intervenção com custo
diário $c=\num{500000}$ dólares. Desprezando impostos, descontos e custos variáveis, o
custo de oportunidade é
\[
      C_{\text{produção}} = q\,p\,d
      = \num{10000}\times\num{70}\times 12
      = \text{US\$}\,\num{8400000}.
\]
O custo direto é $C_{\text{intervenção}}=c\times4=\text{US\$}\,\num{2000000}$, e a
estimativa total é \textbf{US\$\,10,4 milhões}. Se uma ação antecipada reduzisse ambos
os períodos em $30\,\%$, sob a hipótese simplificadora de custos proporcionais ao
tempo, a economia estimada seria $0{,}30\times 10{,}4 =
\textbf{US\$\,3,12 milhões}$. O cálculo é uma análise de cenário, não uma previsão: seu
resultado muda linearmente com $q$, $p$, $d$ e $c$.
\end{exemplo}

\begin{exemplo}{Da janela de detecção ao valor}
Considere o caso do hidrato, discutido em profundidade no \cref{cap:closedworld}.
O intervalo disponível para agir depende do processo e da implantação; a fonte do 3W
consultada nesta revisão não fornece uma faixa temporal que permita fixá-lo. Ainda
assim, o comprimento da janela determina quanto histórico o sistema precisa acumular
antes de produzir uma decisão e deve ser comparado ao intervalo operacional disponível.

Esse é o motivo pelo qual, no \cref{cap:closedworld}, dedicamos uma tabela inteira ao
efeito do comprimento da janela sobre a acurácia. A escolha não é apenas ajuste de
hiperparâmetro: ela combina desempenho observado e restrição operacional, que precisa
ser definida para cada implantação.
\end{exemplo}

Há ainda um custo que não aparece em planilha: o \termo{custo de credibilidade}{alarm
fatigue}. Um sistema que gera falsos alarmes em excesso é desligado --- não
formalmente, mas na prática, quando o operador aprende a ignorá-lo. Esse é o motivo
pelo qual, ao longo do livro, a taxa de falsos positivos recebe tanta atenção quanto a
acurácia, e por que a acurácia balanceada (\cref{cap:aprendizado}) é preferida à
acurácia simples.

\section{Três problemas, não um}
\label{sec:introducao-tresproblemas}

É comum ouvir que ``monitoramento de poços é um problema de classificação''. Não é.
São três problemas, com requisitos diferentes e, como veremos, soluções diferentes.

\begin{description}[leftmargin=0pt,style=nextline,itemsep=0.6em]

\item[Detecção.] \emph{Alguma coisa mudou?} É uma decisão binária sobre uma janela ou
sobre cada amostra de um sinal. Não requer conhecimento de classes; requer apenas um
modelo de normalidade. É o problema atacado por autoencoders com limiar de erro de
reconstrução (\cref{cap:dual}) e, mais tarde, por segmentação supervisionada
(\cref{cap:segmentacao}).

\item[Classificação.] \emph{O que mudou?} Requer um catálogo de classes e, se feita de
forma ingênua, requer que o catálogo seja completo. É o problema atacado pelo ensemble
de classificadores binários (\cref{cap:binarios}) e pela classificação por distância no
espaço latente (\cref{cap:mahalanobis}).

\item[Explicação.] \emph{Por que o sistema acha isso?} Requer que a decisão seja
apresentada em termos que um engenheiro de produção reconheça: qual sensor mudou, em
que direção, com que magnitude, e por que essa combinação é compatível com a hipótese
proposta. É o problema atacado pela camada de agente de linguagem
(\cref{cap:agente,cap:resultados-agente}).

\end{description}

A tese central deste livro é que \textbf{esses três problemas exigem componentes
distintos e que a tentativa de resolvê-los com um único modelo monolítico é a origem
da maior parte das falhas observadas na prática}. Um classificador multiclasse bem
treinado é excelente no segundo problema, péssimo no primeiro quando a anomalia é
inédita, e simplesmente incapaz do terceiro.

\section{A hipótese de mundo fechado}

\begin{definicao}{Hipótese de mundo fechado}
Diz-se que um sistema de classificação opera sob a \termo{hipótese de mundo
fechado}{closed-world assumption} quando pressupõe que toda instância a ser
classificada pertence a uma das classes vistas durante o treinamento. Formalmente, se
$\mathcal{K}$ é o conjunto de classes conhecidas, o sistema modela
$\prob(y \mid \mat{X})$ com suporte restrito a $\mathcal{K}$ e satisfazendo
$\sum_{k \in \mathcal{K}} \prob(y = k \mid \mat{X}) = 1$ para toda entrada $\mat{X}$.
\end{definicao}

A definição parece inocente. A consequência não é. A restrição
$\sum_k \prob(y=k \mid \mat{X}) = 1$ é imposta pela função \ing{softmax} na camada de
saída de praticamente toda rede neural de classificação, e ela significa que \emph{não
existe massa de probabilidade disponível para a hipótese ``nenhuma das anteriores''}.
Quando uma instância de uma classe inédita é apresentada, o modelo é obrigado, por
construção, a distribuir a probabilidade entre as classes que conhece. Pior: não há
nada que force essa distribuição a ser uniforme. É rotineiro observar uma rede
atribuir noventa e cinco por cento de confiança a uma classe errada diante de uma
entrada que ela nunca viu.

\begin{destaque}[A confiança não é a probabilidade]
A saída de uma \ing{softmax} é frequentemente lida como ``a probabilidade de a classe
ser $k$''. Ela não é. Ela é ``a probabilidade de a classe ser $k$, \emph{dado que a
classe está entre as que eu conheço}''. O condicionamento é omitido na leitura
informal, e é exatamente onde mora o erro.
\end{destaque}

O contraponto formal da hipótese de mundo fechado é o \termo{reconhecimento em conjunto
aberto}{open-set recognition}, formulado por \citet{Scheirer2013}, e sua extensão
dinâmica, o \termo{aprendizado de mundo aberto}{open-world learning}, que não apenas
rejeita o desconhecido como o incorpora ao catálogo. Esses conceitos são desenvolvidos
no \cref{cap:novidades} e constituem o eixo da \cref{parte:base}.

\begin{exemplo}{O experimento que dá nome ao livro}
No \cref{cap:closedworld} reproduzimos um experimento simples e decisivo. Treina-se um
sistema de classificação com seis das nove classes de anomalia do conjunto 3W,
escondendo deliberadamente a classe 8 (hidrato na linha de produção). Em seguida,
apresentam-se instâncias dessa classe oculta.

Um classificador multiclasse convencional distribui essas instâncias entre as seis
classes que conhece, com confiança alta e sem qualquer sinal de alerta. O ensemble de
classificadores binários, ao contrário, sinaliza \textbf{89\%} das instâncias como
``nenhuma das classes treinadas''.

A diferença não está na capacidade de representação --- as arquiteturas são
essencialmente as mesmas. Está na estrutura da decisão.
\end{exemplo}

\section{O percurso deste livro}

A narrativa se desenvolve em cinco movimentos.

\paragraph{Parte I --- Fundamentos (\cref{cap:introducao} a \cref{cap:novidades}).}
Construímos o vocabulário. O \cref{cap:dominio} apresenta a arquitetura física de um
poço submarino, seus sensores e a taxonomia de eventos. O \cref{cap:dados} disseca o
conjunto de dados 3W e suas armadilhas. O \cref{cap:aprendizado} estabelece o
arcabouço de avaliação, com ênfase em métricas robustas a desbalanceamento. O
\cref{cap:profundo} percorre as arquiteturas de aprendizado profundo empregadas no
livro. O \cref{cap:novidades} formaliza o problema de mundo aberto e as ferramentas
para atacá-lo.

\paragraph{Parte II --- A base (\cref{cap:dual} a \cref{cap:mundoaberto}).}
Apresentamos o sistema completo desenvolvido na tese de doutorado que originou esta
linha de pesquisa \citep{LopesTese2025}. O \cref{cap:dual} descreve o sistema dual ---
diagrama de decisão mais autoencoder recorrente --- efetivamente implantado em campo.
O \cref{cap:binarios} introduz o ensemble de classificadores binários um-contra-todos e
o valida analiticamente no caso do hidrato. O \cref{cap:mundoaberto} fecha o ciclo com
o pipeline de seis estágios de aprendizado de mundo aberto.

\paragraph{Parte III --- A evolução (\cref{cap:closedworld} a \cref{cap:mahalanobis}).}
Três limitações da base são atacadas. O \cref{cap:closedworld} aprofunda e valida o
ensemble binário \citep{Lopes2026PST}, incluindo a comparação com um modelo
termodinâmico analítico de formação de hidrato. O \cref{cap:segmentacao} substitui a
detecção por erro de reconstrução pela segmentação temporal supervisionada com uma
U-Net unidimensional. O \cref{cap:mahalanobis} substitui a detecção de novidade por
ensemble pela distância de Mahalanobis calculada no espaço latente de um classificador
multiclasse \citep{Lopes2026OTC}.

\paragraph{Parte IV --- A explicação (\cref{cap:llm} a \cref{cap:resultados-agente}).}
Acrescentamos a camada que nenhum dos sistemas anteriores possuía. O \cref{cap:llm}
introduz modelos de linguagem de grande porte com o rigor necessário. O
\cref{cap:agente} descreve a engenharia da camada de agente: extração de descritores
numéricos, construção de perfis de classe e formulação do \ing{prompt}. O
\cref{cap:resultados-agente} apresenta os três estudos que delimitam, com números, o
que o agente sabe e o que ele não sabe fazer.

\paragraph{Parte V --- Síntese (\cref{cap:sintese} a \cref{cap:futuro}).}
O \cref{cap:sintese} consolida todos os resultados em uma tabela única e propõe uma
arquitetura de referência integrada, com divisão de trabalho por classe. O
\cref{cap:operacao} trata de implantação, interface com o operador, deriva de dados e
governança. O \cref{cap:futuro} expõe as limitações remanescentes e as questões em
aberto.

\begin{destaque}[Um livro que mostra as falhas]
Uma advertência de método. Este livro reporta resultados ruins com o mesmo destaque
que reporta resultados bons. A classe 4 (instabilidade de fluxo) atravessa todos os
capítulos com desempenho medíocre; a classe 9 (hidrato na linha de serviço) idem; a
detecção precoce medida por SoftED fica em 0,22 no \cref{cap:segmentacao}; o agente de
linguagem acerta apenas 35,1\% na primeira tentativa no \cref{cap:resultados-agente}.
Esses números estão aqui porque são informativos: eles delimitam a fronteira do que a
metodologia atual alcança, e é nessa fronteira que o trabalho futuro tem valor.
\end{destaque}

\resumocapitulo{%
O monitoramento de poços submarinos é um problema de inferência a partir de telemetria
esparsa, com consequências econômicas mensuráveis em tempo não produtivo. Os eventos
indesejáveis são heterogêneos em escala temporal e em assinatura, o que torna
improvável que um único modelo os cubra bem. Três problemas distintos convivem sob o
rótulo de ``monitoramento'': detectar que algo mudou, classificar o que mudou e
explicar por que o sistema chegou a essa conclusão. A hipótese de mundo fechado ---
pressupor que toda instância pertence a uma classe conhecida --- é uma consequência
estrutural da formulação usual da classificação e produz falhas silenciosas diante de
eventos inéditos. Superá-la é o fio condutor deste livro, que percorre cinco
movimentos: fundamentos, a base metodológica original, sua evolução, a camada de
explicação e a síntese.}

\begin{exercicios}
\item Explique, com suas palavras, por que a restrição
      $\sum_{k \in \mathcal{K}} \prob(y=k \mid \mat{X}) = 1$ imposta pela \ing{softmax}
      impede que um classificador expresse a hipótese ``nenhuma das anteriores''.
      Proponha duas modificações da camada de saída que removeriam essa restrição e
      discuta o que cada uma custaria em termos de treinamento.

\item Um poço produz \num{20000} barris de óleo por dia. Considere um preço de
      \num{70} dólares por barril e um custo diário de intervenção com sonda de
      \num{500000} dólares. Estime o custo total de um evento de hidrato que causa
      \num{12} dias de NPT, dos quais \num{4} exigem sonda. Em seguida, estime a
      economia obtida por um sistema que reduza esse NPT em \SI{30}{\percent}. Discuta
      quais hipóteses do seu cálculo são as mais frágeis.

\item Classifique cada um dos quatro eventos citados na \cref{sec:introducao-oquedaerrado}
      quanto à escala temporal característica (segundos, minutos, horas, meses) e
      quanto ao tipo de assinatura (degrau, rampa, oscilação, perda de sinal).
      Argumente qual tipo de modelo --- recorrente, convolucional ou baseado em regras
      --- você esperaria que funcionasse melhor em cada caso, e registre sua previsão.
      Ao terminar a \cref{parte:evolucao}, confronte-a com os resultados reportados.

\item Discuta a seguinte afirmação: ``um sistema de monitoramento com 99\,\% de acurácia
      é bom o suficiente''. Considere um poço monitorado a cada segundo e uma taxa de
      eventos reais de um por mês. Quantos falsos alarmes por dia esse sistema
      produziria se os 1\,\% de erro se distribuíssem uniformemente? O que isso implica
      sobre a escolha da métrica de avaliação?

\item Considere os três problemas da \cref{sec:introducao-tresproblemas}. Para cada um,
      identifique qual seria a consequência operacional de uma falha: o que acontece se
      a detecção falha? E se a classificação falha, mas a detecção funciona? E se ambas
      funcionam, mas não há explicação?
\end{exercicios}

\begin{leituras}
\item \citet{VARGAS2019} apresentam o conjunto de dados 3W e, mais importante para este
      capítulo, o contexto operacional que motivou sua construção. É a leitura
      obrigatória para quem quiser entender a taxonomia de eventos a partir da prática
      de campo.
\item \citet{Scheirer2013} formalizam o reconhecimento em conjunto aberto. O artigo é
      anterior à popularização do aprendizado profundo e, por isso mesmo, expõe o
      problema com clareza incomum.
\item \citet{Pimentel2014review} oferecem uma revisão ampla de detecção de novidade,
      útil para situar as abordagens deste livro em um panorama maior.
\item \citet{LopesTese2025} é a tese que originou a \cref{parte:base}; os capítulos
      introdutórios contêm um levantamento detalhado do contexto industrial brasileiro.
\end{leituras}
